\section{Lessons from Other Species: Whales, Dolphins, and Finches}
\label{sec:crossspecies}

Dogs are not the first species for which decoding has been attempted, and
the comparison is instructive. The flagship projects differ from the canine
case in one decisive respect: their subjects often possess \emph{discrete,
culturally transmitted vocal units} (whale codas, dolphin signature
whistles, birdsong syllables) that invite combinatorial analysis. Dogs, by
contrast, signal with graded, context-bound affective displays. Reading the
cross-species literature through the ladder of Section~\ref{sec:ladder}
clarifies both what is transferable and what is not.
Figure~\ref{fig:roadmap} places these milestones on a capability timeline;
Table~\ref{tab:milestones} gives the details.

% ===========================================================================
% FIGURE 2 — Capability Roadmap (hero figure)
% ===========================================================================
\begin{figure}[t]
\centering
% Figure 2 — Capability Roadmap (hero figure).
% Compiled standalone via fig2_standalone.tex; also \input by sections/05_crossspecies.tex.
\begin{tikzpicture}[
  font=\sffamily,
  >={Latex[length=2.4mm,width=1.8mm]},
  x=0.40cm,y=1cm,
  tracktitle/.style={font=\sffamily\small\bfseries,anchor=south west},
  projection/.style={font=\sffamily\scriptsize\itshape,anchor=north west,text=ink!78},
  legend/.style={font=\sffamily\scriptsize,anchor=west},
]
  \draw[->,ink!70,line width=0.8pt] (0,0) -- (34.6,0);
  \foreach \yr/\x in {2005/3,2010/8,2015/13,2020/18,2025/23,2030/28,2035/33}{
    \draw[ink!55] (\x,0.07) -- (\x,-0.07);
    \node[font=\sffamily\scriptsize,anchor=north] at (\x,-0.1) {\yr};
  }

  \draw[dashed,rosedeep,line width=0.9pt] (24,-0.35) -- (24,4.55);
  \node[rounded corners=4pt,fill=rosedeep,text=white,inner sep=3pt,
    font=\sffamily\scriptsize\bfseries] at (24,4.8) {NOW $\cdot$ 2026};

  \node[tracktitle] at (0,3.98) {Affect classification \textmd{(Rung 2)}};
  \fill[sky!55,rounded corners=2pt] (0,3.42) rectangle (24,3.94);
  \fill[pattern=north east lines,pattern color=skydeep,rounded corners=2pt]
    (24,3.42) rectangle (33,3.94);
  \draw[ink!55,line width=0.5pt,rounded corners=2pt] (0,3.42) rectangle (33,3.94);
  \node[font=\sffamily\scriptsize,anchor=west] at (0.4,3.68)
    {\textbf{validated} $\cdot$ products ship today};
  \node[projection] at (24.4,3.3) {calibrated welfare monitoring $\sim$2027--30};

  \node[tracktitle] at (14,2.48) {Intent \& context \textmd{(Rung 3)}};
  \fill[sage!55,rounded corners=2pt] (14,1.92) rectangle (24,2.44);
  \fill[pattern=north east lines,pattern color=sagedeep,rounded corners=2pt]
    (24,1.92) rectangle (31,2.44);
  \draw[ink!55,line width=0.5pt,rounded corners=2pt] (14,1.92) rectangle (31,2.44);
  \node[font=\sffamily\scriptsize,anchor=west] at (14.4,2.18)
    {\textbf{emerging} $\cdot$ soundboards + fusion};
  \node[projection] at (24.4,1.8) {validated natural-signal intent $\sim$2030+};

  \node[tracktitle] at (24,0.98) {Open translation \textmd{(Rungs 4--5)}};
  \fill[pattern=north east lines,pattern color=rosedeep,rounded corners=2pt]
    (24,0.42) rectangle (33,0.94);
  \draw[rosedeep,dashed,line width=0.6pt,rounded corners=2pt]
    (24,0.42) rectangle (33,0.94);
  \node[font=\sffamily\scriptsize\itshape,anchor=west,text=rosedeep] at (24.4,0.68)
    {speculative};
  \node[projection] at (24.4,0.36) {no credible timeline};

  \fill[sky!55] (0,-0.98) rectangle (0.6,-0.74);
  \node[legend] at (0.75,-0.86) {validated / shipping};
  \fill[pattern=north east lines,pattern color=ink!60] (8.0,-0.98) rectangle (8.6,-0.74);
  \node[legend] at (8.75,-0.86) {projected (less certain $\rightarrow$)};
  \draw[rosedeep,dashed,line width=0.8pt] (0,-1.34) -- (0.8,-1.34);
  \node[legend] at (0.95,-1.34) {frontier / possibly ill-posed};
\end{tikzpicture}
\caption{\textbf{Capability roadmap, 2002--2035.} Three tracks track the
translation ladder over time. Solid bars denote validated, shipping
capability; hatched bars denote projection, with confidence decreasing to
the right. \textbf{Track A} (affect classification, Rung 2) is mature today
and heads toward calibrated welfare monitoring. \textbf{Track B} (intent and
context, Rung 3) is emerging through soundboards and multimodal fusion, with
validated natural-signal intent plausible after 2030. \textbf{Track C} (open
referential translation and dialogue, Rungs 4--5) has no credible timeline
and may be permanently ill-posed for natural canine communication. The
vertical dashed line marks the present (2026). Method-enabling milestones
(CETI, DolphinGemma, foundation models) are detailed in
Table~\ref{tab:milestones}; they advance the \emph{tools}, not the
\emph{semantics}.}
\label{fig:roadmap}
\end{figure}

\subsection{Project CETI: structure is not meaning}
The Cetacean Translation Initiative is the field's most ambitious program.
Its 2024 \emph{Nature Communications} analysis of 8{,}719 codas from roughly
60 East Caribbean sperm whales demonstrated a context-dependent combinatorial
system---rhythm, tempo, rubato, and ornamentation recombine to expand the
form inventory by nearly an order of magnitude~\cite{sharma2024whale}. The
2022 program paper lays out the interdisciplinary
ambition~\cite{andreas2022ceti}, and the 2025 WhAM model learns to complete
and generate whale-like codas~\cite{wham2025}; a late-2025 preprint even
proposes a translative model of whale
vocalization~\cite{arxiv2025whaletrans}. Yet the investigators are careful:
the ``phonetic alphabet'' describes how signals are \emph{built}, not what
they \emph{mean}. For dogs, the transferable lesson is the methodology
(large corpora, structure discovery, playback validation); the
\emph{non}-transferable part is the assumption of a discrete combinatorial
inventory, which canine graded displays may simply lack.

\subsection{DolphinGemma and CHAT: grounding through interaction}
Google DeepMind's DolphinGemma, with Georgia Tech and the Wild Dolphin
Project, is a roughly 400-million-parameter audio sequence model built on
Gemma with a SoundStream tokenizer, trained on four decades of Atlantic
spotted dolphin recordings and small enough to run on a Pixel
phone~\cite{google2025dolphingemma}. Its companion CHAT device does not try
to translate natural whistles directly; instead it assigns \emph{synthetic}
whistles to coveted objects and lets dolphins request them by name---a
grounding-through-interaction protocol. This is the single most important
design lesson for canine work: meaning is established by \emph{consequence}
(the dolphin gets the sargassum), not by passive correlation. A dog
soundboard is the terrestrial analogue.

\subsection{The Earth Species Project: from encoding to interaction}
The Earth Species Project (ESP), founded in 2018, supplies much of the
field's open infrastructure: the AVES self-supervised
encoder~\cite{hagiwara2023aves}, contributions to NatureLM-audio and the
BEANS benchmark~\cite{robinson2024naturelm,hagiwara2023beans}, and a 2026
zebra-finch model (ZF-AIM) that detects, predicts, and generates calls in
real time, engaging birds in turn-taking exchanges~\cite{esp2026finch}.
ESP's trajectory---encode, then discover structure, then \emph{interact}---%
is the pipeline of Section~\ref{sec:pipeline} made flesh. It also surfaces
the standing critique: Yossi Yovel and others caution that models may
over-rely on acoustic similarity, clustering calls that sound alike but mean
different things~\cite{yovel2025critique}.

\subsection{What transfers to dogs---and what does not}
Three things transfer: the self-supervised encoders, the structure-discovery
toolkit, and above all the interaction-grounding/validation protocol. Three
things do not: the assumption of discrete culturally transmitted units, the
luxury of decades of labeled wild corpora, and the clean playback
accessibility of whales and birds. Dogs are domestic, multimodal, and
graded; their advantage is legibility and proximity, their disadvantage is
the absence of a syntax to decode. The cross-species evidence thus
\emph{supports} the ladder's verdict: the tools for Rungs 1--3 are arriving
fast, while Rungs 4--5 await a discovery about canine semantics that may
never come, because there may be nothing of that kind to discover.

\begin{table}[t]
\centering
\caption{Selected milestones in cross-species communication decoding.
``Type'' distinguishes peer-reviewed structure results (R), generative /
foundation models (M), interactive grounding protocols (I), and consumer
products (P). These advance the toolbox; none yet delivers validated
sentence-level translation for any species.}
\label{tab:milestones}
\small
\begin{tabular}{@{}p{1.0cm}p{4.4cm}p{0.9cm}p{5.7cm}@{}}
\toprule
\textbf{Year} & \textbf{Milestone} & \textbf{Type} & \textbf{Significance} \\
\midrule
2002 & Bow-Lingual~\cite{bowlingual2002} & P & Six-emotion bark matcher; Ig Nobel; the consumer archetype \\
2005 & ELTE bark psychophysics~\cite{pongracz2005barks} & R & Humans classify bark context above chance \\
2013 & DogFACS~\cite{waller2013dogfacs} & R & Objective facial-action coding for dogs \\
2018 & ESP founded; ICASSP bark ML~\cite{icassp2018dogbark} & M/R & Open infrastructure; automated context/affect \\
2022 & CETI program paper~\cite{andreas2022ceti} & R & Roadmap for whale decoding \\
2023 & AVES; BEANS~\cite{hagiwara2023aves,hagiwara2023beans} & M & Self-supervised animal-sound encoder + benchmark \\
2024 & Whale phonetic alphabet~\cite{sharma2024whale} & R & Combinatorial coda structure (form, not meaning) \\
2024 & Soundboard studies~\cite{bastos2024plos,bastos2024scirep} & R/I & Canine word comprehension + intentional combination \\
2024 & NatureLM-audio~\cite{robinson2024naturelm} & M & Audio--language foundation model for bioacoustics \\
2025 & DolphinGemma + CHAT~\cite{google2025dolphingemma} & M/I & Sequence model + object-grounded interaction \\
2025 & WhAM~\cite{wham2025} & M & Generative whale-coda model \\
2026 & ZF-AIM finch interaction~\cite{esp2026finch} & I & Real-time turn-taking with a songbird \\
\bottomrule
\end{tabular}
\end{table}